\section{幻灯片补充页}

\subsection{行为对齐维度}
\clearpage
\noindent\textbf{行为对齐维度}\\
该幻灯片描绘了显式奖励、隐式偏好和安全 guardrail 三层结构之间的依赖关系，对应本文中的 reward signal breakdown。
\begin{center}
\includegraphics[width=0.92\textwidth,page=11]{08_cs224r_reward_learning_2025.pdf}
\end{center}

\subsection{Agent Swarm 流程}
\clearpage
\noindent\textbf{Agent Swarm 流程}\\
该图展示了多 Agent 如何共享 reward trace 以及如何分工处理 metric evaluation，与本章的 Operations Playbook 形成呼应。
\begin{center}
\includegraphics[width=0.92\textwidth,page=12]{08_cs224r_reward_learning_2025.pdf}
\end{center}

\subsection{数据质量控制}
\clearpage
\noindent\textbf{数据质量控制}\\
Data curation pipeline、filter pipes、human annotation loops 的分层结构可用于本章“Batch 数据”章节的参考。
\begin{center}
\includegraphics[width=0.92\textwidth,page=13]{08_cs224r_reward_learning_2025.pdf}
\end{center}

\subsection{治理机制}
\clearpage
\noindent\textbf{治理机制}\\
解释了 audit logs、human review、policy rollback 的连接，对应治理章节中的 checklist。
\begin{center}
\includegraphics[width=0.92\textwidth,page=14]{08_cs224r_reward_learning_2025.pdf}
\end{center}

\subsection{Sim-to-Real 约束}
\clearpage
\noindent\textbf{Sim-to-Real 约束}\\
Slides 强调 domain randomization 与 reward smoothing 的结合，这是我们在“Sim-to-Real”部分补充的内容。
\begin{center}
\includegraphics[width=0.92\textwidth,page=15]{08_cs224r_reward_learning_2025.pdf}
\end{center}

\subsection{Ops Playbook}
\clearpage
\noindent\textbf{Ops Playbook}\\
\begin{center}
\includegraphics[width=0.92\textwidth,page=16]{08_cs224r_reward_learning_2025.pdf}
\end{center}
\noindent 图示展示各 checkpoint 与 runbook，正面支持本章“Observability Playbook”中提出的 checklist。

\subsection{Evidence Matrix}
\clearpage
\noindent\textbf{Evidence Matrix}\\
列出了 reward signal、policy、safety eval 的映射，方便评估不同输入的 evidence quality。
\begin{center}
\includegraphics[width=0.92\textwidth,page=17]{08_cs224r_reward_learning_2025.pdf}
\end{center}

\subsection{未来研究方向}
\clearpage
\noindent\textbf{未来研究方向}\\
列出 CAI、Verifier、Adaptive Reward 等待探索的方向，是本讲“未来方向”部分的直观延伸。
\begin{center}
\includegraphics[width=0.92\textwidth,page=18]{08_cs224r_reward_learning_2025.pdf}
\end{center}
